% ===========================================================================
\section{Methods}
\label{sec:methods}
% ===========================================================================

\subsection{One generation figure, mass-balanced across fates}
\label{sec:partition}

The sector's structure follows from a single decision: municipal solid waste is one national
quantity that must be divided, not three quantities to be estimated separately.

Annual generation is
\begin{equation}
G \;=\; 365 \sum_{d} \left( P^{u}_{d} g_{u} + P^{r}_{d} g_{r} \right),
\end{equation}
over districts $d$, with $g_u$ and $g_r$ the urban and rural per-capita rates. The collected
fraction, which is what reaches a disposal site,
\begin{equation}
D \;=\; 365 \sum_{d} \left( P^{u}_{d} g_{u} (1 - u_{u}) + P^{r}_{d} g_{r} (1 - u_{r}) \right),
\end{equation}
uses the uncollected fractions $u_u$ and $u_r$, and the uncollected remainder is $U = G - D$
by construction rather than by a second estimate. Of that remainder, the burned quantity is
$B = f_b U$.

The identity $D + U = G$ is asserted in code with exact equality, and a test fails if it ever
ceases to hold. (Tier status is stated per component rather than per sector, in
\ref{app:tiers}: the categories share one country-specific activity basis but differ in how
much of each calculation is a default.) This is the property that separate estimation cannot provide. If landfill
tonnage comes from one source and open-burning tonnage from another, their sum is unconstrained
and can exceed or fall short of national generation without anything in the calculation
noticing; here, an argument about the collection rate moves emissions between 4A and 4C2 and
changes the total only through the different methane yields of the two fates.

For FY2023--24 this gives \SI{31.56}{\mega\tonne} generated, \SI{15.88}{\mega\tonne} collected to
disposal, \SI{15.67}{\mega\tonne} uncollected and \SI{9.40}{\mega\tonne} openly burned. The
partition is published as an executable fate ledger
(\texttt{outputs/waste\_ledger\_2024.csv}, the same claimant structure as the transport fuel
ledger), whose writer raises if either tier fails to close and whose final row is the fate
this inventory does not model: \SI{6.27}{\mega\tonne} of uncollected waste that does not
burn---dumped on land or tipped into watercourses, carried explicitly rather than silently.

\subsection{Solid-waste disposal: First Order Decay (4A)}

Methane from a disposal site in year $t$ depends on every year's deposit still decaying, so
IPCC Volume~5 Chapter~3's first-order decay model is run over a deposit history rather than
applied to a single year's tonnage. Decomposable degradable organic carbon accumulates and
decays as
\begin{align}
\mathrm{DDOCm}_{a}(t) &= W(t)\,\mathrm{DOC}\cdot\mathrm{DOC_f}\cdot\mathrm{MCF}
   + \mathrm{DDOCm}_{a}(t{-}1)\,e^{-k}, \\
\mathrm{CH_4}(t) &= \left[\mathrm{DDOCm}_{a}(t{-}1)\left(1 - e^{-k}\right)\right] F
   \cdot \tfrac{16}{12},
\end{align}
with $k = \num{0.065}$, $\mathrm{DOC_f} = \num{0.5}$, $\mathrm{MCF} = \num{0.6}$ and
$F = \num{0.5}$.

Two parameters are deliberately zero. No methane recovery is credited, because none is
documented at national scale and crediting an undocumented recovery would understate the
category. The oxidation factor is zero for the same class of reason: IPCC's non-zero oxidation
applies to managed sites with a cover layer, and Pakistani disposal is predominantly
unmanaged, where the thin intermittent cover that drives oxidation is absent.

\subsubsection{Spin-up, and why it is reported rather than assumed}

The activity panel begins in 2000, but waste deposited before then is still emitting. Starting
the model at the panel's first year with an empty site would assert that Pakistan landfilled
nothing before 2000, which is false and biases the estimate low.

The deposit series is therefore back-projected to 1960 at \SI{40}{\percent} of the panel's
first-year tonnage, and the model is run from there. Because that back-projection is an
assumption rather than data, the sector reports what it is worth: without spin-up the category
would be \SI{8.07}{\mega\tonne} instead of \SI{9.17}{\mega\tonne}, so pre-2000 deposits
contribute \SI{12}{\percent} of today's landfill methane. A reader who prefers to start at the
panel can read the alternative directly from the output file rather than recomputing it.

\subsubsection{The monthly state engine, and why it is one recursion}

The national release resolves the disposal stream into fiscal months, and the
engine behind those months is this model run at a monthly step rather than a
reshaped annual answer: one continuous single-pool recursion from the 1960
spin-up, $S_m = S_{m-1}\,e^{-k\,d_m/365.25} + W_m$, where $d_m$ is the month's
exact calendar length (leap Februaries included) and $W_m$ deposits the month's
own calendar year's tonnage, days-weighted within the year. Running one recursion
rather than a per-year replay matters for three named reasons: state flows across
fiscal boundaries (June's stock is July's opening balance, so no seam appears at
the year change); each month deposits its \emph{own} year's waste (a per-year
replay had injected the panel's maximum year into every year's stock); and the
reaction delay between deposit and the onset of decay is a declared parameter
(\texttt{swds\_delay\_months}, the IPCC six-month convention) applied as a shape
refinement only. The single pool is kept deliberately --- the annual model carries
one decay constant, and a cohort structure would claim a waste-composition split
no Pakistani record supports. The annual control is preserved exactly by
construction: each fiscal year's twelve shares are normalised against its own
annual figure, so the monthly engine decides only \emph{when} within the year,
never how much, and the within-year spread it computes is printed into the
month rows' basis rather than asserted in prose.

\subsection{Open burning (4C2)}

IPCC Equation~5.1 applied to the burned quantity $B$, taking that chapter's own
allowance to combine the dry-matter, carbon-content and fossil-fraction terms into one
coefficient rather than carrying them separately:
\begin{equation}
\mathrm{CO_2} \;=\; B \cdot \mathrm{CF_{fossil}} \cdot \mathrm{OF}
  \cdot \tfrac{44}{12},
\end{equation}
where $\mathrm{CF_{fossil}}$ is the fossil carbon per \emph{wet} tonne given in
Section~\ref{sec:data} and the oxidation factor is $\mathrm{OF} = \num{0.71}$. Dry matter
does not appear here because it is already inside that coefficient; it enters this category
only through the \NtwoO{} factor, which IPCC states per dry tonne. The burned quantity $B$
is itself IPCC Equation~5.7.

Only the fossil carbon fraction produces reportable \COtwo{}. The biogenic fraction is
reported as a memo item at \SI{3520}{\giga\gram} \COtwo{}, outside the national total, because
that carbon was fixed from the atmosphere within the year. Methane and nitrous oxide from the
same burning \emph{are} counted regardless of carbon origin: incomplete combustion is a net
source however the carbon entered the waste stream.

\subsection{Domestic wastewater (4D1)}

Tier~1 over the sanitation-resolved population. Total organically degradable material is
\begin{equation}
\mathrm{TOW} \;=\; P \cdot \mathrm{BOD} \cdot 365 \cdot I,
\end{equation}
with $I = \num{1.25}$ the industrial-and-commercial correction, and methane is
\begin{equation}
\mathrm{CH_4} \;=\; \sum_{j} \mathrm{TOW}_{j} \cdot B_{0} \cdot \mathrm{MCF}_{j},
\end{equation}
summed over sanitation pathways $j$ with $B_0 = \SI{0.6}{\kilo\gram}$ \CHfour{} per kilogram
BOD. The pathway correction factors are the subject of Section~\ref{sec:discussion}.

Nitrous oxide is computed from protein consumption rather than from BOD, by the method of
the 2019 Refinement (Volume~5 Chapter~6, Equations~6.7 to 6.10a) \citep{ipcc2019}. Protein
supply is reduced to protein consumed (a factor of \num{0.96}) and its nitrogen (\num{0.16})
is allocated to the same sanitation pathways the methane uses, with factors of \num{1.1} for
household products, \num{1.1} for non-consumed protein (the chapter's updated Table~6.11,
taken over its text, which gives no increment where food waste goes with solid waste) and,
for industrial and commercial co-discharge, \num{1.25} on the sewered population only.
Treatment removes part of each pathway's nitrogen before discharge---\num{0.15} in septic
tanks, \num{0.12} in latrines and \num{0.40} in the treated share of sewered flow, which is
declared at \SI{1}{\percent} and not measured---and the nitrogen entering a treatment plant
emits there at \num{0.016}. What is left is summed over every pathway and emits at the
effluent factor of \num{0.005}; the whole population stays in the estimate. The result is
\SI{9.71}{\giga\gram} of \NtwoO{}. The single equation of the 2006 Guidelines, which this
inventory used before and which treats one nitrogen total on protein supply as effluent,
is retained as the alternative and gives \SI{11.45}{\giga\gram}. Neither method reads BOD or
the sewer correction factor. This is why nitrous oxide does not move under either of the
sensitivities reported below, and why a sensitivity row that reported methane alone would be
misread as a category total---a distinction this sector's output file previously got wrong
and now states explicitly.
